\subsection{The case that X/H is paracompact}

If X/H is paracompact, we shall first see that the vector spaces \(\mathcal{K}^{\chi}(X)\) , for variable \(\chi\) , are all isomorphic to each other, and in particular isomorphic to \(\mathcal{K}^{1}(X)\) .

PROPOSITION 7. --- Assume X/H paracompact. Let \(\chi\) be a continuous representation of H in \(\mathbf{R}_{+}^{*}\) .

\begin{enumerate}
\def\labelenumi{\alph{enumi})}
\item
  There exists on \(\mathbf{X}\) a continuous function \(r\) , with values \(>0\) , such that \(r(x\xi) = \chi (\xi)r(x)\) for all \(x\in \mathbf{X}\) and \(\xi \in \mathbf{H}\) .
\item
  The mapping \(g \mapsto g / r\) is an isomorphism of the vector space \(\mathcal{H}^{\chi}(\mathbf{X})\) onto the vector space \(\mathcal{H}^{1}(\mathbf{X})\) .
\end{enumerate}

Let us apply Prop. 1 of No.~1 on taking f to be a function \(\geqslant0\) that is not identically zero on any orbit (this is possible by Lemma 1 of Appendix 1); then \(r = f^{x}\) satisfies the properties of a). The assertion b) is obvious.

PROPOSITION 8. --- Assume X/H paracompact. There exists a continuous function \(h \geqslant 0\) on X, whose support has compact intersection with the saturation of every compact subset of X, and is such that \(h^{b}=1\) . For such a function, one has \(g=\left(h\cdot(g\circ\pi)\right)^{b}\) for every continuous function g on X/H.

Let us apply Prop. 1 of No.~1, with \(\chi = 1\) , on taking for \(f\) a function \(\geqslant 0\) that is not identically zero on any orbit. We have \(f^1(x) > 0\) at every point \(x\) of X. Set \(h = f / f^1\) . Then \(h^1 = f^1 / f^1 = 1\) , therefore \(h^b = 1\) . It follows that if \(g\) is a continuous function on X/H, then \(\left(h \cdot (g \circ \pi)\right)^b = h^b \cdot g = g\) .

Remarks. --- 1) In particular, let X be a locally compact space on which a discrete group D operates continuously and properly on the right; suppose X/D is paracompact. Then, there exists a continuous function \(h \geqslant 0\) on X whose support has compact intersection with the saturation of every compact subset of X, and is such that \(\sum_{d \in D} h(x d) = 1\) for every \(x \in X\) (all terms of the sum being zero except for a finite number of them).

\begin{enumerate}
\def\labelenumi{\arabic{enumi})}
\setcounter{enumi}{1}
\tightlist
\item
  Let us conserve the hypotheses and notations of Prop. 8. The mapping \(g \mapsto h \cdot (g \circ \pi)\) is a continuous mapping of \(\mathcal{K}(\mathrm{X / H})\) into \(\mathcal{K}(\mathrm{X})\) that is a right inverse of the mapping \(f \mapsto f^{\flat}\) . Consequently, every bounded (resp. compact) subset of \(\mathcal{K}(\mathrm{X / H})\) is the image of a bounded (resp. compact) subset of \(\mathcal{K}(\mathrm{X})\) . From this, one deduces immediately that the mapping \(\lambda \mapsto \lambda^{\sharp}\) is again an isomorphism of \(\mathcal{M}(\mathrm{X / H})\) onto a closed linear subspace of \(\mathcal{M}(\mathrm{X})\) when these spaces are equipped with the topology of bounded (resp. compact) convergence.
\end{enumerate}

PROPOSITION 9. --- We conserve the hypotheses and notations of Prop. 8. Let \(\lambda\) be a positive measure on X/H.

\begin{enumerate}
\def\labelenumi{\alph{enumi})}
\item
  The pair \((\pi, h)\) is \(\lambda^{\sharp}\) -adapted, and \(\int_{\mathbf{X}} h(x) \varepsilon_{\pi(x)} d\lambda^{\sharp}(x) = \lambda\) .
\item
  The mapping \(\pi\) is proper for the measure \(h\cdot \lambda^{\sharp}\) , and \(\pi (h\cdot \lambda^{\sharp}) = \lambda\) .
\item
  Let \(\mathbf{k}\) be a function on \(\mathrm{X / H}\) , with values in a Banach space or in \(\overline{\mathbf{R}}\) . For \(\mathbf{k}\) to be measurable (resp. locally integrable, essentially integrable, integrable) for \(\lambda\) , it is necessary and sufficient that \(h\cdot (\mathbf{k}\circ \pi)\) be so for \(\lambda^{\sharp}\) ; and, if \(\mathbf{k}\) is essentially integrable for \(\lambda\) , then
\end{enumerate}

\[
\int_ {\mathrm{X} / \mathrm{H}} \mathbf {k} d \lambda = \int_ {\mathrm{X}} h \cdot (\mathbf {k} \circ \pi) d \lambda^ {\sharp}.\tag{14}
\]

Let \(f \in \mathcal{K}(\mathrm{X}/\mathrm{H})\) . Then \(h \cdot (f \circ \pi) \in \mathcal{K}(\mathrm{X})\) and

\[
\int_ {\mathrm{X}} h (x) f (\pi (x)) d \lambda^ {\sharp} (x) = \int_ {\mathrm{X} / \mathrm{H}} f (\dot {x}) d \lambda (\dot {x}) \int_ {\mathrm{H}} h (x \xi) d \beta (\xi) = \int_ {\mathrm{X} / \mathrm{H}} f (\dot {x}) d \lambda (\dot {x}),
\]

whence a). The assertion b) is proved similarly. The assertions of c) concerning measurability, essential integrability and formula (14) may then be obtained by applying the results of Ch. V (§4, Prop. 3, §5, Th. 1, §4, Th. 2). If k is λ-integrable, then \(h \cdot (\mathbf{k} \circ \pi)\) is \(\lambda^{\sharp}\) -integrable (Ch. V, §3, No.~3, Th. 1). If \(h \cdot (\mathbf{k} \circ \pi)\) is \(\lambda^{\sharp}\) -integrable, Prop. 5 proves that \((h \cdot (\mathbf{k} \circ \pi))^b = h^b \cdot \mathbf{k} = \mathbf{k}\) is λ-integrable. If k is locally λ-integrable, then \(h \cdot (\mathbf{k} \circ \pi)\) is locally \(\lambda^{\sharp}\) -integrable (Prop. 6). Finally, suppose \(h \cdot (\mathbf{k} \circ \pi)\) locally \(\lambda^{\sharp}\) -integrable; for every \(f \in \mathcal{H}(\mathrm{X}/\mathrm{H})\) , \(h \cdot (\mathbf{k} \circ \pi) \cdot (f \circ \pi)\) has compact support, and

\[
\left| h \cdot (\mathbf {k} \circ \pi) \cdot (f \circ \pi) \right| \leqslant \mathrm{M} \left| h \cdot (\mathbf {k} \circ \pi) \right|,
\]

where \(M = \sup |f|\) ; therefore \(h \cdot ((\mathbf{k}f) \circ \pi)\) is \(\lambda^{\sharp}\) -integrable, consequently kf is \(\lambda\) -integrable, by what has already been proved; this proves that k is indeed locally \(\lambda\) -integrable.

COROLLARY. --- The continuous linear mapping \(f \mapsto f^{b}\) of \(\mathrm{L}^{1}(\mathrm{X}, \lambda^{\sharp})\) into \(\mathrm{L}^{1}(\mathrm{X}/\mathrm{H}, \lambda)\) defined by Prop. 5 is surjective.

Suppose first that X/H is paracompact and let h be a function on X satisfying the conditions of Prop. 8. If k is a \(\lambda\) -integrable numerical function on X/H, then \(h \cdot (k \circ \pi)\) is \(\lambda^{\sharp}\) -integrable and \((h \cdot (k \circ \pi))^{\flat} = k\) (Prop. 9).

In the general case, let \(u \in \mathrm{L}^1(\mathrm{X/H},\lambda)\) . There exists a function \(f \in \mathcal{L}^1(\mathrm{X/H},\lambda)\) with class \(u\) and zero outside a countable union of compact sets \(\mathbf{K}_n\) . Let us define recursively a sequence of relatively compact open sets \(\mathrm{U}_n\) of \(\mathrm{X/H}\) such that \(\mathrm{U}_{n+1} \supset \mathrm{K}_n \cup \overline{\mathrm{U}}_n\) , and let \(V\) be the union of the \(\mathrm{U}_n\) . Then \(V\) is an open subset of \(\mathrm{X/H}\) , a countable union of compact subsets \(\overline{\mathrm{U}}_n\) , hence is paracompact (GT, I, §9, No.~10, Th. 5). Set \(Y = \overline{\pi}^1(V)\) and let \(\lambda_V\) (resp. \(\lambda_Y^\sharp\) ) be the measure induced by \(\lambda\) (resp. \(\lambda^\sharp\) ) on \(V\) (resp. \(Y\) ). It is clear that \(Y/H\) may be identified with \(V\) (GT, I, §3, Prop. 10) and that \(\lambda_Y^\sharp\) may be identified with \((\lambda_V)^{\sharp}\) . Moreover, \(f\) is zero outside \(V\) and belongs to \(\mathcal{L}^1(V,\lambda_V)\) . Therefore, there exists \(g \in \mathcal{L}^1(Y,\lambda_Y^\sharp)\) such that \(g^b = f\) almost everywhere in \(V\) . Extending \(g\) by 0 on \(X - Y\) , one obtains a function \(g_1 \in \mathcal{L}^1(X,\lambda^\sharp)\) , and it is clear that the class of \(g_1^b\) in \(L^1(X/H,\lambda)\) is none other than \(u\) .

Remark 3. --- Suppose X/H paracompact, and let us retain the notations of Proposition 9. The mapping \(k \mapsto h \cdot (k \circ \pi)\) of \(L^1(X/H, \lambda)\) into \(L^1(X, \lambda^\sharp)\) is then isometric by (14) and is a right inverse of the mapping \(f \mapsto f^\flat\) of \(L^1(X, \lambda^\sharp)\) onto \(L^1(X/H, \lambda)\) .

